\subsection{An arithmetic source for the exponents \texorpdfstring{$5\cdot3^r$}{5 times powers of three}}
\label{subsec:lm-fifteen-source}

Let $A=\operatorname{Jac}(C)$, where $C$ is the genus-two curve
$v^2+v=u^5$, and put
\[
 F=\mathbb Q(\zeta_5),\qquad K=\mathbb Q(A[3]),\qquad S=\{3,5\}.
\]
The following calculation controls the full arithmetic action, not only
one rational isogeny factor of a higher Fermat Jacobian.

\begin{proposition}[The residual field]
\label{prop:lm-a5-residual}
The field $K$ contains $F$ and $\mathbb Q(\zeta_3)$, and
\[
 [K:\mathbb Q]=80,\qquad
 \operatorname{Gal}(K/F)\simeq C_{20}.
\]
There is exactly one prime of $K$ above $3$ and exactly one above $5$.
Their ramification indices and residue degrees are respectively
$(20,4)$ and $(40,2)$.  If $H\simeq C_5$ is the unique subgroup of
order five in $\operatorname{Gal}(K/F)$ and $K_0=K^H$, then $K_0/\mathbb Q$
is Galois of degree $16$, and
\begin{equation}
 |D_K|=3^{76}5^{94},\qquad |D_{K_0}|=3^{12}5^{14}.
 \label{eq:lm-residual-discriminants}
\end{equation}
The ramification index of $K_0$ above $5$ is $8$.
\end{proposition}
\begin{proof}
The automorphism $(u,v)\mapsto(\zeta_5u,v)$ gives the full CM ring
$\mathbb Z[\zeta_5]$ on $A_F$. Its CM type is primitive, since the
cyclic quartic field $F$ has no imaginary quadratic subfield; hence
$A$ is absolutely simple.  Since $3$ is inert in $F$,
$A[3]$ is a one-dimensional vector space over
$\mathbb Z[\zeta_5]/3\simeq\mathbb F_{81}$.
If a Galois element acts trivially on $A[3]$, its conjugation on
the reduction of $\zeta_5$ is trivial.  The element $\zeta_5$
has exact order five in $\mathbb F_{81}^{\times}$; hence the
element fixes $F$.  This proves $F\subset K$.
The Weil pairing gives $\mathbb Q(\zeta_3)\subset K$.

The principal polarization has Rosati involution equal to
complex conjugation on $F$.  On $\mathbb F_{81}$ this is
$a\mapsto a^9$.  Consequently the image of $G_F$ is contained in
\[
 \{a\in\mathbb F_{81}^{\times}:a^{10}\in\mathbb F_3^{\times}\}
 =\mu_{20}.
\]
This is a direct upper bound; no general index assertion for CM
representations is needed.

The curve has good reduction at $3$.  The differential basis
$du/(2v+1),u\,du/(2v+1)$ gives CM type $\{1,2\}$ on $F$,
whose reflex type is $\{1,3\}$.  The local CM reciprocity
formula identifies the inertia action at $3$ with the inverse
reflex norm on $\mathcal O_{F_3}^{\times}$.
Since Frobenius acts on $\zeta_5$ by exponent $3$, its reduction is
\[
 \mathbb F_{81}^{\times}\longrightarrow\mathbb F_{81}^{\times},
 \qquad a\longmapsto a^{-4}.
\]
Its image has order $80/\gcd(80,4)=20$.  Thus the global
image over $F$ is exactly $C_{20}$, and the local inertia at
$3$ already fills it.  The field $F_3/\mathbb Q_3$ is
unramified of degree four.  It follows that the decomposition
group at $3$ is the entire group of order $80$, with
ramification index $20$ and residue degree four.
The reciprocity formula used here is stated in
\citet[Theorem~5.1, pp.~483--484]{Lombardo2017}; the
finite unit factor is trivial at a place of good reduction.

For the other ramified prime, use the isomorphism
\[
 x=4u,\qquad y=16(2v+1),\qquad y^2=x^5+256.
\]
\citet[Proposition~2.2(4)--(5), author p.~4]{StollYang2003}
give no conductor at $2$ and, up to inverse and coefficient
conjugation, the following character on
units at $\lambda=1-\zeta_5$:
\[
 \alpha\longmapsto
 \left(\frac{\alpha}{\lambda}\right)_2
 \rho(\alpha)^{q(256)},\qquad
 q(256)=\frac{256^4-1}{5}\equiv4\pmod5.
\]
Here $\rho$ has exact order five and conductor $\lambda^2$;
the quadratic factor has conductor $\lambda$.
Both factors retain their orders modulo $3$.
Thus inertia in $K/F$ at $5$ is $C_{10}$, with wild part
$C_5$ and tame part $C_2$.  The inertia index over
$\mathbb Q_5$ is therefore $4\cdot10=40$.
Because $K$ contains $\zeta_3$, its residue degree at $5$
is at least two.  The product already equals $[K:\mathbb Q]$,
so the prime above $5$ is unique and its residue degree is two.
There is no ramification elsewhere, by good reduction of
$A$ away from $5$ and the prime-to-residue-characteristic
torsion criterion.

Apply the conductor--discriminant formula to the abelian
extension $K/F$ of degree twenty.  At $3$, all nineteen
nontrivial characters have conductor exponent one; the
norm of the prime $(3)$ of $F$ is $3^4$.
At $5$, sixteen characters are nontrivial on the wild
$C_5$ and have conductor exponent two, two further
characters are nontrivial only on the tame $C_2$ and
have exponent one, and one nontrivial character is
unramified.  Thus
\[
 |D_K|=(5^3)^{20}\,3^{4\cdot19}\,5^{16\cdot2+2}
       =3^{76}5^{94}.
\]
The subgroup $H$ is characteristic in $C_{20}$ and hence
normal in $\operatorname{Gal}(K/\mathbb Q)$.
The extension $K_0/F$ has group $C_4$.
Its three nontrivial characters have exponent one at
$3$; at $5$ two have exponent one and one is unramified.
Consequently
\[
 |D_{K_0}|=(5^3)^4\,3^{4\cdot3}\,5^2=3^{12}5^{14}.
\]
Its ramification index at $5$ is $4\cdot2=8$.
\end{proof}

\begin{lemma}[The exact residual group and ramification]
\label{lem:lm-quintic-residual-group}
With the notation of Proposition~\ref{prop:lm-a5-residual},
\[
 \operatorname{Gal}(K/\mathbb Q)=
 \langle\tau,\varphi\mid\tau^{20}=1,\quad
 \varphi\tau\varphi^{-1}=\tau^3,\quad \varphi^4=\tau^{10}\rangle.
\]
Thus its extension of $C_{20}$ by $C_4$ is nonsplit: every lift of the
indicated generator of $C_4$ has order eight. At five the inertia group
is $C_5\rtimes C_8$, with the indicated generator of $C_8$ acting by
third powers on $C_5$. Its lower ramification groups are
\[
 |G_0|=40,\qquad G_1=G_2=C_5,\qquad G_j=1\quad(j\ge3).
\]
\end{lemma}
\begin{proof}
Put $k=\mathbb F_{81}$ and identify $A[3]$ with $k$.
The Rosati adjoint relation writes the alternating Weil pairing as
\[
 \langle x,y\rangle=\operatorname{Tr}_{k/\mathbb F_3}(cxy^9),
 \qquad c\ne0,\quad c^9=-c.
\]
Indeed every linear functional on $k$ is a trace functional, and the
adjoint of a scalar $a$ is $a^9$; alternation gives the condition on $c$.
Every residual Galois element is semilinear, of the form
$x\mapsto ax^{3^i}$. A lift of the third-power generator of
$\operatorname{Gal}(\mathbb Q(\zeta_5)/\mathbb Q)$ has $i=1$.
For its multiplier $m\in\mathbb F_3^\times$, the trace pairing gives
\[
 (ca^{10})^{27}=mc,\qquad a^{10}=mc^2.
\]
Consequently its fourth power is the scalar
\[
 a^{1+3+9+27}=a^{40}=(mc^2)^4=c^8=-1.
\]
For a generator $\tau$ of the scalar subgroup $C_{20}$, this scalar
is $\tau^{10}$, and conjugation acts on $\tau$ by its third power.
The displayed presentation has at most eighty normal forms
$\tau^a\varphi^i$, $0\le a<20$, $0\le i<4$. The already proved degree
is eighty, so it is the exact presentation. Multiplying a lift by a
scalar does not change its fourth power, since $1+3+9+27=40$.

At five the mod-three cyclotomic multiplier is unramified. The known
inertia group of order forty is therefore its kernel. It surjects
onto $C_4$, since the fifth cyclotomic field is totally ramified at
five. A lift in inertia of the indicated generator has order eight,
and together with its normal wild subgroup $C_5$ it generates
$C_5\rtimes C_8$ of order forty.

Finally the unique prime above five has residue degree two. The
exponent ninety-four in the established discriminant therefore makes
the local different exponent forty-seven. All positive lower groups
are subgroups of the wild group $C_5$. If exactly $t$ of them are
nontrivial, the different formula is
$47=(40-1)+t(5-1)$, whence $t=2$. This proves the asserted filtration.
\end{proof}

\begin{theorem}[Local arithmetic source dominance]
\label{thm:lm-fifteen-source}
Let $K_S(3)$ be the maximal pro-$3$ extension of $K$
unramified outside the primes above $3$ and $5$, and put
\[
 \Gamma=\operatorname{Gal}(K_S(3)/\mathbb Q).
\]
For the unique prime $\mathfrak p\mid3$ of $K$, the maps
\[
 G_{K_{\mathfrak p}}\longrightarrow
 \operatorname{Gal}(K_S(3)/K),\qquad
 G_{\mathbb Q_3}\longrightarrow\Gamma
\]
are surjective.
\end{theorem}
\begin{proof}
Set $G=\operatorname{Gal}(K/\mathbb Q)$ and
\[
 V=\ker\bigl(H^1(G_{K,S},\mathbb F_3)
       \longrightarrow H^1(K_{\mathfrak p},\mathbb F_3)\bigr).
\]
The unique prime above $3$ makes $V$ a $G$-module.  This
module is finite dimensional, and is semisimple because
$3\nmid|G|=80$.

First suppose $V^H\ne0$, where $H=C_5$ is as above.
Inflation--restriction for the degree-five extension
$K/K_0$ identifies $H^1(G_{K_0,S},\mathbb F_3)$ with
$H^1(G_{K,S},\mathbb F_3)^H$.  The same restriction is
injective locally, since $[K_{\mathfrak p}:(K_0)_{\mathfrak p_0}]=5$.
Thus a nonzero element gives a cyclic cubic extension
$L_0/K_0$ unramified outside $S$, and split completely
at the prime above $3$.
Only the primes above $5$ can ramify in $L_0/K_0$;
this ramification is tame and has index at most three.
The tame different formula, and $e_5(K_0/\mathbb Q)=8$,
give
\[
 \operatorname{rd}(L_0)
 \le 3^{12/16}5^{14/16}\,5^{(1-1/3)/8}
 =3^{3/4}5^{23/24}<10.659.
\]
But $[L_0:\mathbb Q]=48$, whereas Odlyzko's unconditional
bound at degree at least $48$ is $\operatorname{rd}>13.736$.
This is a contradiction, so $V^H=0$.

If $V\ne0$, choose an irreducible $G$-submodule $W\subset V$.
Every nontrivial irreducible $\mathbb F_3[C_5]$-module has
dimension four: $\Phi_5(T)$ is irreducible over $\mathbb F_3$
because $3$ has order four modulo five.
Since $W^H=0$, we have $\dim W\ge4$.
The common kernel of the characters in $W$ cuts out an
elementary abelian extension $L/K$ of degree
$3^{\dim W}\ge81$.  It is Galois over $\mathbb Q$ because
$W$ is $G$-stable.  It is split at $3$ and unramified
outside $S$.

At $5$, tame inertia in an elementary abelian $3$-extension
is cyclic of exponent three, so its order is at most three.
Using $e_5(K/\mathbb Q)=40$ and
\eqref{eq:lm-residual-discriminants} gives
\[
 \operatorname{rd}(L)
 \le 3^{76/80}5^{94/80}\,5^{(1-1/3)/40}
 =3^{19/20}5^{143/120}<19.329.
\]
On the other hand $[L:\mathbb Q]\ge80\cdot81=6480$.
Odlyzko's unconditional bound at every degree at least
$4800$ is $\operatorname{rd}>21.843$.
This contradiction proves $V=0$.

For precision, these are the degree-$48$ and degree-$4800$ rows
of \citet[Table~2]{Odlyzko1976}. The accompanying description,
paragraphs (i)--(ii), specifies that the complex-column
bound applies to every number field of degree at least
the stated threshold.  No generalized Riemann hypothesis
or numerical class-group assumption enters the argument.

Let $P=\operatorname{Gal}(K_S(3)/K)$ and let $D$ be the
closed image of $G_{K_{\mathfrak p}}$ in $P$.
If $D$ were proper, a finite $3$-group quotient of $P$
would have proper image of $D$; that image lies in an
index-three maximal subgroup.  The resulting nonzero
character $P\to\mathbb F_3$ would vanish on $D$,
contradicting $V=0$.  Thus $D=P$.
Finally $G_{\mathbb Q_3}$ already surjects onto
$\operatorname{Gal}(K/\mathbb Q)$ by
Proposition~\ref{prop:lm-a5-residual}, and its intersection
with $P$ contains all of $P$ by what was just proved.
Hence it surjects onto $\Gamma$.
\end{proof}

\begin{lemma}[The full geometric action over the residual field]
\label{lem:lm-five-power-action}
Let $n=5\cdot3^r$, $r\ge0$, and let $N_n$ be the geometric
pro-$3$ fundamental group of the affine Fermat curve $U_n$.
With a standard unit tangent at any rational coordinate cusp,
the actual $G_{\mathbb Q}$-action on $N_n$ is unramified
outside $S$, factors through $\Gamma$, and has pro-$3$
image on $G_K$.
For every finite free $\mathbb Z_3$-module $M$ obtained
from this action by tensor operations and subquotients,
\[
 H^1(G_{\mathbb Q,S},M)\longrightarrow H^1(\mathbb Q_3,M)
 \quad\text{is injective}.
\]
\end{lemma}
\begin{proof}
First consider $n=5$.  The established three cyclic maps
give a $\mathbb Q$-isogeny $J_5\sim A^3$.
The image of $G_K$ on $T_3A$ lies in the kernel of
$\operatorname{GL}(T_3A)\to\operatorname{GL}(A[3])$,
which is pro-$3$.  The rational Tate-module isomorphism
coming from the isogeny therefore identifies the image
on $V_3J_5$ with a pro-$3$ group.  Being pro-$3$ is a
topological property and is unaffected by changing a
stable lattice.  In particular the image on $J_5[3]$
is a finite $3$-group, so all its composition factors
over $G_K$ are trivial.  This argument does not require
the isogeny degree to be prime to three.

The boundary of $X_5$ is three copies of $\mu_5$.
Because $K$ contains both $\mu_5$ and $\mu_3$, all
composition factors of the boundary module modulo
$3$ are also trivial over $G_K$.
The exact boundary--compact sequence for $N_5^{\mathrm{ab}}$
therefore shows that $G_K$ has a finite $3$-group image
on $N_5/\Phi(N_5)$: a representation all of whose
composition factors are trivial can be put in upper
unipotent form.
The kernel of
$\operatorname{Aut}(N_5)\to\operatorname{Aut}(N_5/\Phi(N_5))$
is pro-$3$ for a finitely generated free pro-$3$ group.
Consequently the actual image on $N_5$ is pro-$3$.

For $n=5\cdot3^r$, the coordinate-power map $U_n\to U_5$
has geometrically Galois group $\mu_{3^r}^2$.
Thus $N_n$ identifies with an open subgroup of $N_5$.
At the rational cusp $(0:-1:1)$ its local parameters
satisfy $u_5=u_n^{3^r}$, so the unit tangents are
compatible and the upstairs action is the restriction
of the downstairs action.  The subgroup is stable for
this action because the map and the compatible tangents
are defined over $\mathbb Q$.
The same statement holds at the other rational cusps
by coordinate permutation.  This is a genuine
$3$-power cover; no prime-to-$3$ cover is treated as an
injection of pro-$3$ groups.

The smooth pair and the chosen extending tangent give
unramified actual action at every prime outside $S$.
Its restriction to $G_K$ has pro-$3$ image, hence
factors through $\operatorname{Gal}(K_S(3)/K)$; the
whole action consequently factors through $\Gamma$.
These properties pass to $M$.

Let $z$ represent a class in $H^1(G_{\mathbb Q,S},M)$.
Over $G_K$ its graph takes values in the pro-$3$ group
$M\rtimes\rho(G_K)$, so both the action and the cocycle
factor through $\Gamma$.
If the local class is zero, subtract one coboundary
so that $z$ is identically zero on $G_{\mathbb Q_3}$.
Theorem~\ref{thm:lm-fifteen-source} makes this cocycle
zero on $\Gamma$, proving injectivity.
\end{proof}

\begin{theorem}[Unconditional mixed exponents $5\cdot3^r$]
\label{thm:lm-five-power-selmer}
For every $r\ge1$ and $n=5\cdot3^r$, the compact
localization criterion $(\mathrm C_n)$ holds:
\[
 \ker\!\left(\operatorname{Sel}(T_3J_n)
       \longrightarrow H^1(\mathbb Q_3,T_3J_n)\right)=0.
\]
Every Selmer section of the geometric pro-$3$ abelian
quotient of $U_n/\mathbb Q$ is globally cuspidal.
The same holds for the full geometric pro-$3$ group,
every finite lower-central quotient, and every finite-derived
quotient (in particular the metabelian quotient).
More generally it holds for any arithmetic-stable quotient
retaining $N_n^{\mathrm{ab}}$ whose ordinary lower-central
factors are torsion-free over $\mathbb Z_3$.
In each case the Selmer set consists of exactly the
three disjoint rational cusp packets parametrized by
$H^1(\mathbb Q,\mathbb Z_3(1))$.
As throughout, a local Selmer representative means the
image of a full point section or a full rational-cusp section;
no global full lift is assumed for a quotient section.
\end{theorem}
\begin{proof}
A compact Selmer class is unramified outside $S$:
there $J_n$ has good reduction and its local Kummer
classes are unramified.  Apply
Lemma~\ref{lem:lm-five-power-action} to $T_3J_n$, a
quotient of $N_n^{\mathrm{ab}}$, to obtain $(\mathrm C_n)$.
The mixed compact criterion then gives the abelian
conclusion.  To spell out its two essential geometric
steps, the morphism $X_n\to X_3$ and the established
cubic compact Selmer input force the local center at
$3$ to be one rational coordinate cusp $P$.
At any $p\notin S$, equality of compact Kummer classes
forces every boundary reduction or cusp center to be
$P$: the difference of distinct coordinate cusps has
exact order $n$, whose $3$-primary part is nonzero.
This is where $r\ge1$ is used.
The boundary component transverse to $I_P$ is therefore
unramified outside $S$ and zero at $3$, so it vanishes
by the same cohomological injection.
The surviving component gives the unique global
cuspidal parameter $\beta$.

For the nonabelian conclusion write
$L_m=\gamma_mR/\gamma_{m+1}R$.
These finite free grades are subquotients of tensor powers of
$N_n^{\mathrm{ab}}$, so their actions factor through $\Gamma$ and
Lemma~\ref{lem:lm-five-power-action} makes their $S$-ramified first
cohomology localize injectively at three. At a good prime, the
compact and boundary factors have strictly positive Frobenius weights,
as do their tensor subquotients. Hence $L_m$ has no global invariants.
Surjectivity of $G_{\mathbb Q_3}\to\Gamma$ gives
$L_m^{G_{\mathbb Q_3}}=0$ as well. We have verified every condition
$(\mathrm G_m)$ in Proposition~\ref{prop:lm-real-mixed-higher-criterion}.
That proposition therefore gives the claimed classification directly
in $R$, using full local representatives and no global lift.
The induction allows arbitrary ramification at five and arbitrary,
possibly infinite, support of the cusp parameter. Only each successive
graded difference is required to be unramified outside $S$.
The free-group and completion lemmas give torsion-free grades for
all the specific quotients listed in the statement.
This proves the $5\cdot3^r$ case of
Theorem~\ref{thm:unrestrictedthreefamiliesintro}.
\end{proof}

\begin{remark}[Scope]\label{rem:lm-five-power-scope}
This proves a mixed-exponent family, including $n=15$,
without computing the Mordell--Weil rank or Tate--Shafarevich
group of $J_{15}$.  It does not assert the same result
for $n=5$, for $5^s3^r$ with $s>1$, or for other
prime-to-three factors.  For $n=5$, the residual/source
calculation remains true but distinct compact cusp
differences have prime-to-three order, so the indicated
boundary-center step is unavailable.  For $s>1$, the
map to $U_5$ has a prime-to-three part and does not
give the open subgroup used in the action lemma.
\end{remark}
